\section{Observational Verification}
\label{sec:verification}

The framework makes quantitative predictions for both Saturn and Jupiter.
We test them against Cassini and Juno observations, treating the two
planets as independent case studies of the same geometric principle.

\subsection{A generalized observable: \texorpdfstring{$\sigma_{\mathrm{geom}}$}{sigma-geom}}
\label{sec:sigma-geom}

To compare polygonal patterns across planets and dynamical mechanisms,
we define a universal observable.  Given $N$ features (cyclone centers,
jet meander vertices) at positions $z_k$ in the complex plane (via
stereographic projection from the pole), the geometric M\"obius parameter
is
\begin{equation}\label{eq:sigma-geom}
  \sigma_{\mathrm{geom}} = \frac{1}{N}\sum_{k=1}^{N}
  \bigl|\!\ln|z_k / z_k^{\mathrm{ideal}}|\bigr|,
\end{equation}
where $z_k^{\mathrm{ideal}}$ are the vertices of the best-fit regular
$N$-gon (optimized over rotation angle).  By construction,
$\sigma_{\mathrm{geom}} = 0$ for a perfect polygon ($|\lam| = 1$) and
$\sigma_{\mathrm{geom}} > 0$ for any distortion.  The observable is
rotation-invariant, scale-invariant, and independent of $N$.

\subsection{Saturn: the hexagonal jet meander}

\subsubsection{Three-epoch data}

Saturn's hexagon has been observed across three epochs spanning
three decades (\cref{tab:epochs}).

\begin{table}[htbp]
\centering
\caption{Saturn hexagon drift rates across three observational epochs.
All angular velocities relative to System~III.}
\label{tab:epochs}
\begin{tabular}{lcccl}
\toprule
Epoch & $\omega_{\mathrm{hex}}$ ($^\circ$/day) & $\omega_{\mathrm{NPS}}$ ($^\circ$/day) & NPS & Source \\
\midrule
1980--81  & $-0.060 \pm 0.014$ & $-0.044 \pm 0.010$ & Yes   & \citet{Godfrey1988} \\
1990--91  & $-0.001 \pm 0.007$ & ---                  & Weak  & \citet{SanchezLavega1993} \\
2008--14  & $-0.036 \pm 0.010$ & ---                  & No    & \citet{SanchezLavega2014} \\
\bottomrule
\end{tabular}
\end{table}

\subsubsection{Wavenumber and amplitude}

The Cassini/ISS zonal wind profile \citep{Antunano2015}, digitized near
the hexagonal jet ($70$--$82^\circ$\,N), gives a peak speed
$U_{\mathrm{obs}} = 120.1$~m/s at $76.6^\circ$\,N with half-width
$2.86^\circ$.  The Rossby stationarity condition $U^* = \beta/k^2$ with
$k = 6/R_{\mathrm{hex}}$ predicts $U^* = 105$~m/s, selecting $n^* = 5.62
\to 6$ as the stationary wavenumber (12\% discrepancy attributable to the
barotropic approximation).

The hexagon meander amplitude, measured from the vertex latitude span
($75^\circ$--$78^\circ$\,N), is $\varepsilon_{\mathrm{obs}} = 0.112$.
The Theorem~5 prediction with the derived matching constant $C = 0.797$
gives $\varepsilon_{\mathrm{pred}} = 0.112$.  The numerical coincidence
with $\varepsilon_{\mathrm{obs}} = 0.112$ should not be overstated:
propagating the 9.8\% Gaussian profile approximation error gives
$\varepsilon_{\mathrm{pred}} = 0.11 \pm 0.11$, so the match is within
one standard deviation of the model uncertainty, not sub-percent.

\subsubsection{NPS coupling and near-resonance}

The North Polar Spot (NPS), a compact anticyclonic vortex near the
hexagon, was prominent during 1980--81 but absent by 2008--14.
The linear coupling model $\omega_{\mathrm{hex}} = \omega_0 +
\gamma\,\omega_{\mathrm{NPS}}$ gives $\gamma = 0.54 \pm 0.41$
(two parameters fitted to two data points with no residual degrees of
freedom; the 1990--91 epoch is not used because $\omega_{\mathrm{NPS}}$
was not measured).  The implied near-resonance period is
\[
  T_{\mathrm{rel}} = \frac{360^\circ}{|\omega_{\mathrm{hex}} -
  \omega_{\mathrm{NPS}}|} = 62 \pm 68\;\text{years},
\]
where the large uncertainty reflects the propagated error bars on both
drift rates.  The multi-decadal timescale is \emph{suggestive} of an
NPS--hexagon coupling but cannot be considered well-constrained by the
available data.

\subsection{Jupiter: the cyclone rings}

Jupiter provides the critical universality test.  Its polar cyclone
rings arise through vortex self-organization, not Rossby wave stationarity
---the Rossby prediction gives $n^* \approx 0.3$ for both poles,
completely failing to match the observed $N = 8$ (north) or $N = 5$
(south).  If the M\"obius framework applies only to Rossby-mediated
polygons, it should fail for Jupiter.  It does not.

\subsubsection{Juno cyclone positions}

Cyclone positions digitized from \citet{Adriani2018} are projected
stereographically from the pole and measured via
$\sigma_{\mathrm{geom}}$ (\cref{eq:sigma-geom}).

\subsubsection{Thomson stability}

The north polar ring ($N = 8$) is unstable without a central vortex
(Havelock: $N \ge 8$ unstable), but Jupiter has a strong central cyclone
at each pole that provides stabilization.  The south polar ring ($N = 5$)
is stable without a center (Havelock: $N \le 7$), with Thomson critical
ratio $\kappa_{\mathrm{crit}} = -1/2$.

\subsection{Cross-planetary comparison}

\paragraph{Quantitative criterion.}
We say a configuration is ``near $|\lam| = 1$'' if its M\"obius parameter
satisfies $\sigma < 1$ (equivalently $r = e^\sigma < e \approx 2.72$).
This threshold is motivated physically: $\sigma < 1$ means the spiral
pitch angle $\psi = \arctan(\alpha/\sigma)$ exceeds $45^\circ$---the
orbit is closer to a circle than to a radial ray.  We further distinguish
\emph{close} ($\sigma < 0.1$, $r < 1.11$) from \emph{moderate}
($0.1 < \sigma < 1$).

For Saturn, we note that two different observables give different values:
\begin{itemize}
  \item $\sigma_{\mathrm{flow}} \approx 0.10$ (from the dynamical BVP
    chain, \cref{sec:theorems})---\emph{close} to $|\lam| = 1$.
  \item $\sigma_{\mathrm{geom}} = 0.50$ (from the vertex distortion
    measure, \cref{sec:sigma-geom})---\emph{moderate}.
\end{itemize}
The difference arises because $\sigma_{\mathrm{geom}}$ measures the
full geometric distortion of the hexagonal meander (amplitude
$\varepsilon = 0.112$), while $\sigma_{\mathrm{flow}}$ measures the
Rossby wave's departure from exact stationarity ($\delta U/U^* = 0.14$).
Both place Saturn within the $\sigma < 1$ criterion; Jupiter sits in
the ``close'' regime by either measure.

\begin{table}[htbp]
\centering
\caption{All three planetary configurations tested against the framework.}
\label{tab:cross-planet}
\begin{tabular}{lccccl}
\toprule
System & $N$ & $\sigma_{\mathrm{geom}}$ & $\sigma_{\mathrm{flow}}$ & Regime & Entry mechanism \\
\midrule
Saturn hexagon    & 6 & 0.50 & 0.10 & Close (flow) / Moderate (geom) & Rossby \\
Jupiter north     & 8 & 0.06 & --- & Close & Thomson + center \\
Jupiter south     & 5 & 0.05 & --- & Close & Thomson \\
\bottomrule
\end{tabular}
\end{table}

Jupiter's cyclone rings ($\sigma_{\mathrm{geom}} < 0.07$) sit close to
$|\lam| = 1$.  Saturn's hexagon ($\sigma_{\mathrm{geom}} = 0.50$,
$r = 1.65$) departs substantially further.  This difference is physically
expected: $\sigma_{\mathrm{geom}}$ measures the \emph{geometric}
distortion of the observed feature from a perfect $N$-gon, and Saturn's
hexagon is a jet meander with amplitude $\varepsilon = 0.112$---a
large-scale deformation of a circular jet that necessarily distorts the
polygon.  Jupiter's cyclone rings, by contrast, are discrete point-like
vortices whose positions can closely approximate a regular polygon even
when the surrounding flow is not circular.  The two cases probe different
regimes of the same framework: Saturn tests the Rossby-mediated
\emph{wave} regime ($\sigma_{\mathrm{flow}} \approx 0.10$ from the
dynamical chain, closer to $|\lam| = 1$ than $\sigma_{\mathrm{geom}}$
suggests); Jupiter tests the vortex-cluster regime where
$\sigma_{\mathrm{geom}}$ directly measures the ring geometry.

A quantitative criterion for ``near $|\lam| = 1$'' depends on the
observable used: for the dynamical M\"obius parameter
$\sigma_{\mathrm{flow}}$, Saturn sits at $\sim 0.10$ (close); for the
geometric observable $\sigma_{\mathrm{geom}}$, Saturn sits at $0.50$
(moderate).  We report both and do not claim tighter agreement than the
data support.

The Rossby mechanism (Theorem~3) fails entirely for Jupiter.  The Thomson
mechanism provides the alternative entry point.  Different dynamics,
same geometric destination.

\begin{table}[htbp]
\centering
\caption{Summary of verified predictions.}
\label{tab:verification}
\begin{tabular}{clccl}
\toprule
\# & Prediction & Value & Observed & Status \\
\midrule
1 & Saturn wavenumber $n^*$ & $5.62$ & $6$ & Correct integer \\
2 & Saturn amplitude $\varepsilon$ & $0.11 \pm 0.11$ & $0.112$ (Cassini) & Consistent \\
3 & Stationary jet speed $U^*$ & $105$~m/s & $120$~m/s & $12\%$ (barotropic) \\
4 & Thomson $\kappa_{\mathrm{crit}}$ ($N=6$) & $-1/4$ & --- & Exact algebraic \\
5 & Branch disconnection & $\sigma = 0$ (within WKB) & 40-year persistence & Proven (WKB) \\
6 & Jupiter $\sigma_{\mathrm{geom}}$ & $< 0.07$ & Near-polygonal & Confirmed \\
\bottomrule
\end{tabular}
\end{table}
